\begin{exercise}{7.2}{4}
    Gegeben ein verallgemeinerter Kreis $M$ und zwei Punkte $p, q$ gibt es
    stets einen verallgemeinerten Kreis $L$, der durch unsere beiden Punkte geht und
    auf $M$ senkrecht steht. Unter der Annahme $p \notin {q, s_M(q)}$ ist L eindeutig bestimmt
\end{exercise}

\textit{Hinweis:} Diese Lösung ist sehr an Soergels Lösung angelehnt, die allgemein ja immer etwas flapsiger sind. Eine detailliertere Lösung wäre gut.
\bigskip

Möbistransformationen sind transitiv auf verallgemeinerten Kreisen und erhalten Orthogonalität. Daher sei o.B.d.A $M = \hat{g}$ eine erweiterte Gerade.
Fallunterscheidungen:
\begin{enumerate}
    \item $p=q$: Offensichtlich gibt es unendlich viele Lösungen.
    \item $p \neq q$ und $p,q \neq \infty$: Betrachte die Mittelsenkrechte $h$ auf der Strecke $[p,q]$.
    \begin{enumerate}
        \item $h$ nicht parallel oder gleich $g$, gibt es einen Schnittpunkt $c$ mit $g$. Der Kreis $L$ durch $p$ mit Zentrum $c$ geht auch durch $q$ und steht senkrecht auf $M$. Diese Lösung ist eindeutig.
        \item Ist $h=g$, so kann $L$ genau so konstruiert werden. Es ist jedoch keine eindeutige Lösung.
        \item Ist $h$ parallel zu $g$, so steht die affine Gerade $\ell$, welche durch $p$ und $q$ geht, senkrecht auf $g$. Durch die Erweiterung mit $\infty$ zu $\widehat{\ell}$ erhalten wir $L$ als eindeutige Lösung.
    \end{enumerate}
    \item $p=\infty$ oder $q=\infty$: Sei o.B.d.A. $p = \infty$. Nehme als $\ell$ die zu $p$ senkrechte Gerade und erweitere sie mit $\infty$, um $\hat{\ell} =: L$ zu erhalten. Diese Lösung ist eindeutig.
\end{enumerate}